\documentclass[12pt]{article}
\usepackage[utf8]{inputenc}
\usepackage[T1]{fontenc}
\usepackage[spanish]{babel}
\usepackage{geometry}
\usepackage{array}
\usepackage{tabularx}
\usepackage{enumitem}
\usepackage{longtable}
\usepackage{amsmath}
\usepackage{fancyhdr}
\usepackage{lastpage}
\usepackage{xcolor}
\geometry{letterpaper, margin=2.2cm}
\pagestyle{fancy}
\fancyhf{}
\lhead{Curso de Programaci\'on}
\rhead{Unidad 4 -- Clave docente}
\cfoot{P\'agina \thepage\ de \pageref{LastPage}}
\setlength{\parindent}{0pt}
\renewcommand{\arraystretch}{1.2}
\begin{document}

\begin{center}
{\LARGE \textbf{Clave docente -- Eval\_U4}}\\[2mm]
{\large Hoja de respuestas y orientaciones de calificaci\'on}\\[1mm]
\end{center}

\section*{Criterios generales de calificaci\'on}
\begin{itemize}[leftmargin=1.2em]
    \item Acepte sin\'onimos t\'ecnicamente correctos en preguntas abiertas, siempre que no contradigan el contenido del curso.
    \item En preguntas de c\'odigo, otorgue puntaje parcial si la l\'ogica es correcta aunque exista un error menor de sintaxis.
    \item Penalice el uso de comprensiones, funciones no vistas o respuestas que rompan la restricci\'on pedag\'ogica de la unidad.
\end{itemize}

\section*{Pregunta 1. Emparejamiento de conceptos (12 puntos)}
Respuestas correctas:
\begin{itemize}[leftmargin=1.2em]
    \item 1. Lista $\rightarrow$ b
    \item 2. Tupla $\rightarrow$ c
    \item 3. Set $\rightarrow$ d
    \item 4. Diccionario $\rightarrow$ a
    \item 5. Mutabilidad $\rightarrow$ e
    \item 6. \texttt{append()} $\rightarrow$ f
\end{itemize}
Puntaje sugerido: 2 puntos por relaci\'on correcta.

\section*{Pregunta 2. Completar frase con palabra clave (12 puntos)}
\begin{enumerate}[leftmargin=1.3em]
    \item mutable
    \item inmutable
    \item key
    \item value
    \item index
    \item set
\end{enumerate}
Puntaje sugerido: 2 puntos por respuesta correcta.

\section*{Pregunta 3. Completar fragmentos de c\'odigo (20 puntos)}
\subsection*{3.1 Registro en listas}
Respuesta esperada:
\begin{verbatim}
temperaturas = [18, 21]
temperaturas.append(24)
print(temperaturas[-1])
\end{verbatim}
Criterio: 4 puntos por \texttt{append(24)} y 4 puntos por \texttt{-1}.

\subsection*{3.2 B\'usqueda en diccionario}
Respuesta esperada:
\begin{verbatim}
print(aeronave["horas"])
\end{verbatim}
Aceptar tambi\'en \texttt{'horas'}.

\subsection*{3.3 Recorrido de una tupla}
Respuesta esperada:
\begin{verbatim}
for modelo in modelos:
    print(modelo)
\end{verbatim}
Puntaje: 6 puntos completos si usa la tupla correcta.

\section*{Pregunta 4. Completar c\'odigo para resolver un problema (18 puntos)}
Respuesta esperada:
\begin{verbatim}
if pieza["horas"] >= pieza["limite"]:
\end{verbatim}
Explicaci\'on esperada: la condici\'on verifica si las horas de uso de la pieza alcanzan o superan el l\'imite establecido, por lo que debe recibir mantenimiento.

Puntaje sugerido:
\begin{itemize}[leftmargin=1.2em]
    \item 10 puntos por la condici\'on correcta.
    \item 8 puntos por la explicaci\'on clara del criterio de comparaci\'on.
\end{itemize}

\section*{Pregunta 5. Deducci\'on de funcionamiento de un fragmento (18 puntos)}
Fragmento final:
\begin{verbatim}
['fuselaje', 'motores', 'tren de aterrizaje']
\end{verbatim}
Respuestas esperadas:
\begin{enumerate}[leftmargin=1.3em]
    \item La lista final es \texttt{['fuselaje', 'motores', 'tren de aterrizaje']}. 
    \item \texttt{append()} agrega \texttt{"tren de aterrizaje"} al final de la lista.
    \item \texttt{remove()} elimina la primera ocurrencia de \texttt{"alas"}.
\end{enumerate}
Puntaje sugerido: 6 puntos por cada inciso.

\section*{Pregunta 6. Auditor\'ia de c\'odigo (20 puntos)}
Respuestas esperadas:
\begin{enumerate}[leftmargin=1.3em]
    \item \texttt{flota} es un diccionario. Cada clave es una matr\'icula y cada valor es una lista con tres datos: modelo, horas y un tercer valor num\'erico que en el c\'odigo se compara como l\'imite o referencia.
    \item Salida esperada:
\begin{verbatim}
HK-101 cumple
HK-202 no cumple
\end{verbatim}
    \item Mejora posible: usar nombres m\'as descriptivos, por ejemplo almacenar cada aeronave como diccionario con llaves \texttt{"modelo"}, \texttt{"horas"} y \texttt{"limite"}; tambi\'en puede extraerse la lista a una variable temporal para evitar indexaci\'on repetida.
\end{enumerate}
Puntaje sugerido:
\begin{itemize}[leftmargin=1.2em]
    \item 6 puntos por identificar la estructura y el contenido.
    \item 6 puntos por la salida correcta.
    \item 8 puntos por una mejora razonable y bien justificada.
\end{itemize}

\section*{Pregunta 7. Completar c\'odigo con estructuras de datos correctas (10 puntos)}
Respuesta esperada:
\begin{verbatim}
inventario = {}
\end{verbatim}
Justificaci\'on esperada: un diccionario es adecuado porque permite asociar cada repuesto con su cantidad mediante una clave \'unica.

Puntaje sugerido:
\begin{itemize}[leftmargin=1.2em]
    \item 4 puntos por \texttt{\{\}}.
    \item 6 puntos por la justificaci\'on conceptual.
\end{itemize}

\section*{Observaciones pedag\'ogicas}
\begin{itemize}[leftmargin=1.2em]
    \item La evaluaci\'on se mantiene dentro de los contenidos observados en la unidad y se alinea con las actividades pr\'acticas del repositorio.
    \item El tiempo estimado de 100 minutos es adecuado para resolver el examen con lectura, an\'alisis y redacci\'on breve.
\end{itemize}

\end{document}
